\section{Strategische Ma{\ss}nahmen}
\label{sec:strategie}

Dieser Teil berichtet ausschlie{\ss}lich, was die Berichte sagen; beurteilt wird in den
Teilen~\ref{sec:optionen} bis~\ref{sec:votum}.

\subsection{Kapazit\"atsausbau in Rotterdam}

Der Ausbau der Raffinerie Rotterdam ist das gr\"o{\ss}te laufende Projekt. Im
Gesch\"aftsbericht 2024 teilte Neste mit, dass sich wegen des angespannten Markts f\"ur
Bauunternehmen der Beginn des kommerziellen Betriebs von 2026 auf 2027 verschiebt und die
Investitionssumme von \je{\RotterdamKostenAlt}{Mrd.\,EUR} auf \je{\RotterdamKostenNeu}{Mrd.\,EUR}
steigt\QL{gbviernein}{GBXXIV}{91}, also um \RotterdamMehrkosten. Nach Abschluss soll die Kapazit\"at f\"ur
erneuerbare Produkte \je{\KapazitaetZiel}{Mio.\,t} betragen\Q{GBXXIV}{8}, gegen\"uber rund
\je{\Kapazitaet}{Mio.\,t} heute\Q{GBXXV}{14}; das sind \je{\KapazitaetPlus}{Mio.\,t} oder
\KapazitaetPlusPz{} mehr. 2025 wurden f\"ur das Projekt Bauzeitzinsen von
\Mio{\BauzinsenAktiviert}{EUR} aktiviert\Q{GBXXV}{210}. Im Halbjahresbericht 2026 hei{\ss}t es,
der Ausbau \enquote{continues to progress}; ein neuer Termin oder eine neue Kostenangabe
wird dort nicht genannt\Q{HJ}{2}.

Zugleich k\"undigte Neste an, die Investitionen au{\ss}erhalb Rotterdams in den folgenden
Jahren auf j\"ahrlich rund \je{\InvestNachRotterdam}{Mrd.\,EUR} zu begrenzen, mit Schwerpunkt
auf Sicherheit und Verl\"asslichkeit\QR{gbviernein}.

\subsection{Leistungsprogramm}

Nach dem Wechsel an der Spitze startete Neste Ende 2024 ein Programm zur Verbesserung der
Ertragskraft. Ziel ist eine Verbesserung des EBITDA um \Mio{\ProgrammZiel}{EUR} Laufrate bis
Ende 2026 gegen\"uber der Basis 2024; Ende 2025 waren \Mio{\ProgrammLaufrate}{EUR}
erreicht\Q{GBXXV}{80}, Ende Juni 2026 \Mio{\ProgrammLaufrateHj}{EUR}. Im zweiten Quartal
2026 trug das Programm \Mio{\ProgrammQuartal}{EUR} zum EBITDA bei\Q{HJ}{1}.

\subsection{Stillst\"ande im zweiten Halbjahr 2026}

F\"ur die zweite Jahresh\"alfte 2026 sind drei planm\"a{\ss}ige Stillst\"ande angek\"undigt:
Porvoo f\"ur rund \StillstandPorvoo{} Wochen ab August, Rotterdam f\"ur rund
\StillstandRotterdam{} Wochen im vierten Quartal und eine Produktionslinie in Singapur f\"ur
rund \StillstandSingapur{} Wochen ab Dezember\QL{hjzwei}{HJ}{2}. F\"ur den Porvoo-Stillstand hat Neste
vorab Produkte hergestellt und eingelagert; der Aufbau des Umlaufverm\"ogens um
\Mio{\QzNwcAufbau}{EUR} im zweiten Quartal geht vor allem darauf zur\"uck\QR{hjzwei}.

\subsection{Markt und Regulierung, wie das Unternehmen sie beschreibt}

Der Vorstandsvorsitzende f\"uhrt das Rekordergebnis des zweiten Quartals auf den Konflikt
im Nahen Osten zur\"uck, der die \"Ol- und Produktm\"arkte \"uber den gr\"o{\ss}ten Teil des
Zeitraums beherrscht und ein \enquote{exceptional market environment} geschaffen habe\QL{hjeins}{HJ}{1}.
Als dauerhafte Treiber nennt er h\"ohere Beimischungspflichten: Die Umsetzung der
Erneuerbare-Energien-Richtlinie RED~III in Deutschland soll die Nachfrage nach erneuerbarem
Diesel 2026 um rund \je{\RedDreiNachfrage}{Mio.\,t} j\"ahrlich erh\"ohen; in den USA hoben
h\"ohere Pflichtmengen die Preise der Handelszertifikate (RIN)\QR{hjeins}. F\"ur den
Luftverkehr schreibt die EU-Verordnung ReFuelEU Aviation einen Anteil nachhaltigen
Flugkraftstoffs von \Pz{\SafQuoteZiel} f\"ur 2025 und \Pz{\SafQuoteZielDreissig} f\"ur 2030
vor; nach Angaben der EASA wurden 2025 \Pz{\SafQuoteIst} erreicht\QW{EASA, Pressemitteilung
zu ReFuelEU Aviation}{quelle:easa}{27.09.2026}.

\subsection{Wirkung auf EBITDA und EBIT}

Renewable Products steigerte das vergleichbare EBITDA 2025 auf \Mio{\RpEbitda}{EUR} nach
\Mio{\RpEbitdaVj}{EUR}, bei einem Absatz von \je{\RpAbsatz}{Mio.\,t} nach
\je{\RpAbsatzVj}{Mio.\,t}\Q{GBXXV}{79}. Im zweiten Quartal 2026 erreichte das Segment
\Mio{\QzRpEbitda}{EUR} nach \Mio{\QzRpEbitdaVj}{EUR}; Oil Products kam auf
\Mio{\QzOpEbitda}{EUR} bei einer Raffineriemarge von \je{\QzOpMarge}{USD/bbl} nach
\je{\QzOpMargeVj}{USD/bbl}\QR{hjeins}. Die Verbesserung von 2025 tr\"agt vor allem die Menge
in Renewable Products; die von 2026 tragen die Margen beider Hauptsegmente.

\subsection{Ausblick gegen Ist}

F\"ur 2026 erwartet Neste einen Absatz von Renewable Products etwa auf dem Niveau von 2025,
einen geringeren Absatz von Oil Products wegen des Stillstands und Investitionen ohne M\&A von
rund \je{\InvestPrognoseHj}{Mrd.\,EUR}\QR{hjzwei}. Im Gesch\"aftsbericht 2025 lautete die
Investitionsprognose noch \je{\ProgXXVIVon}{} bis \je{\ProgXXVIBis}{Mrd.\,EUR}\Q{GBXXV}{85};
sie liegt jetzt am oberen Rand. Eine Ergebnisprognose in Euro gibt Neste nicht.

Die einzige Gr\"o{\ss}e, die Neste jedes Jahr beziffert, sind die Investitionen. Gegen die
Prognose wie zuerst gegeben:

\prognosetabelle

In \PrognoseUnter{} von f\"unf Jahren lagen die Investitionen unter der zuerst genannten
Zahl. F\"ur 2021 senkte Neste die Prognose unterj\"ahrig von \je{\ProgXXI}{} auf
\je{\ProgXXIAngepasst}{Mrd.\,EUR}\Q{GBXXI}{152}. Zur Verkaufsmarge von Renewable Products
gab Neste 2024 eine Spanne vor und senkte sie im Mai und im September zweimal\Q{GBXXIV}{98}.
